\begin{table}[!t]
\caption{Accuracy--Latency Frontier at Level 12 ($n{=}200$)}
\label{tab:frontier}
\centering
\begin{tabular}{@{}rrrr@{}}
\toprule
\textbf{Budget (s)} & \textbf{Mean MASE} & \textbf{Mean CV RMSE} & \textbf{Mean cost (s)} \\
\midrule
0.01      & 0.748 & 12.544 & 0.0002 \\
0.10      & 0.745 & 12.513 & 0.004 \\
0.50      & 0.745 & 12.006 & 0.277 \\
1.00      & 0.716 & 11.839 & 0.519 \\
\textbf{2.00} & \textbf{0.698} & 11.724 & 0.925 \\
5.00      & 0.717 & 11.112 & 1.980 \\
10.00     & 0.733 & 10.781 & 3.383 \\
30.00     & 0.736 & 10.576 & 7.358 \\
$\infty$  & 0.735 & 10.562 & 10.843 \\
\bottomrule
\end{tabular}
\vspace{2pt}
\begin{minipage}{\columnwidth}
\footnotesize
Mean CV RMSE---the quantity the selector minimizes---decreases monotonically as
the budget relaxes, confirming the selection rule behaves correctly. Realized
MASE does not: it reaches a minimum near 2~seconds and degrades thereafter.
Admitting the models with the lowest cross-validated error worsens out-of-sample
accuracy, so the budget operates as regularization against validation-fold
overfitting. The minimum uses 8.5\% of the computation required by unconstrained
selection while achieving 5\% lower error.
\end{minipage}
\end{table}
